% Generated by Claude Opus 5 (Anthropic) — 2026-09-13
% Final campaign overhead_aw_20260913_025626 (fishwait11, optimized overlay,
% per_instance + kernel sched_switch). Model counts: fish_20260913_014334 (LTTng+nsys #1).

Out of the profiled session FISH generated an application model with 88 executors, 241 nodes and 3451 callbacks linked by 6199 interactions on the callback level.

\begin{table}[h]
\centering
\footnotesize
\caption{Overhead added with profiling on Autoware, the values taken from Autoware's own nodes reporting on \texttt{/diagnostics}, in the steady state of a 30\,s
rosbag replay: mean over three runs (upper line) with the delta against the baseline, and their range (lower line). Processing time is the mean over the nodes that report one, concatenation counted over complete point-cloud sets. The row below end-to-end latency restricts it to the chain of nodes that are active, and measured, during the rosbag replay. The last two rows are
  the diagnostics that would report a failure.}
\label{tab:overhead-aw}
\begin{tabular}{@{}lrrr@{}}
\hline
 & baseline & LTTng & LTTng+nsys \\ \hline
\begin{tabular}[c]{@{}l@{}}CPU nodes, 17 of 239\\ processing time [ms]\end{tabular}
  & \begin{tabular}[c]{@{}r@{}}2.79\\ \scriptsize 2.76--2.84\end{tabular}
  & \begin{tabular}[c]{@{}r@{}}3.06 (+9.5\%)\\ \scriptsize 2.96--3.17\end{tabular}
  & \begin{tabular}[c]{@{}r@{}}3.09 (+10.6\%)\\ \scriptsize 3.03--3.21\end{tabular} \\ \hline
\begin{tabular}[c]{@{}l@{}}GPU nodes, 2 of 2\\ processing time [ms]\end{tabular}
  & \begin{tabular}[c]{@{}r@{}}13.95\\ \scriptsize 13.72--14.08\end{tabular}
  & \begin{tabular}[c]{@{}r@{}}13.99 (+0.3\%)\\ \scriptsize 13.91--14.08\end{tabular}
  & \begin{tabular}[c]{@{}r@{}}14.21 (+1.9\%)\\ \scriptsize 13.91--14.45\end{tabular} \\ \hline
\begin{tabular}[c]{@{}l@{}}end-to-end\\ latency [ms]\end{tabular}
  & \begin{tabular}[c]{@{}r@{}}412\\ \scriptsize 402--418\end{tabular}
  & \begin{tabular}[c]{@{}r@{}}419 (+1.7\%)\\ \scriptsize 416--421\end{tabular}
  & \begin{tabular}[c]{@{}r@{}}423 (+2.8\%)\\ \scriptsize 413--443\end{tabular} \\ \hline
\begin{tabular}[c]{@{}l@{}}active chain\\ latency [ms]\end{tabular}
  & \begin{tabular}[c]{@{}r@{}}154\\ \scriptsize 144--161\end{tabular}
  & \begin{tabular}[c]{@{}r@{}}161 (+4.6\%)\\ \scriptsize 158--163\end{tabular}
  & \begin{tabular}[c]{@{}r@{}}166 (+7.6\%)\\ \scriptsize 155--185\end{tabular} \\ \hline
\begin{tabular}[c]{@{}l@{}}\textbf{centerpoint missed}\\ \textbf{cycles}\end{tabular} & \textbf{0} & \textbf{0} & \textbf{0} \\ \hline
\begin{tabular}[c]{@{}l@{}}\textbf{topic state monitor}\\ \textbf{violations}\end{tabular} & \textbf{0} & \textbf{0} & \textbf{0} \\ \hline
\end{tabular}
\end{table}

\noindent
\textbf{Overhead added with profiling:}
The untraced baseline carries no trace, so the only per-node values comparable across the three modes are the processing times Autoware's nodes report on \texttt{/diagnostics}: 19 of the 241 do, both GPU nodes and the 17 CPU nodes of the lidar preprocessing chain, NDT and prediction, and Table~\ref{tab:overhead-aw} averages them. End-to-end latency grows by 1.7\,\% under LTTng and 2.8\,\% with GPU profiling on top. In our experiments we find that tracing changes how long the application takes, not what it does: the perception pipeline keeps producing detected objects at the same rate (9.9\,Hz in every mode), no centerpoint cycle is missed, and no topic state monitor leaves
\textsc{ok} in any of the nine runs, meaning the added latency stays below what
Autoware's own supervision treats as a fault.
The GPU nodes are the least affected, and one of them, the occupancy grid, even gets faster (3.64 to 3.38\,ms): its callback ends in a blocking device-to-host copy that waits behind the centerpoint inference sharing the GPU, so the time it reports measures how the two callbacks collide on the accelerator, which tracing perturbs, rather than the cost of observing it.
